% Regression: planes_keys_test — exercises frame and side-policy keys across the documented plane geometries.
% Formula: geometry-only; no tensor identity.
% planes_keys_test — acceptance for the frame / side-policy
% key work (sweep-verdict worklist items 1-7).  Every picture cites the
% worklist item and the sweep measurement that motivated it.  Expected
% warnings (and ONLY these): planes-interleave + sheet-coincide (item 1,
% part 3),
% plane-rise-low + plane-slant-band (item 2, part 6), plane-tall-window
% (item 3, part 8) -- plus plane-tall-window at part 5: the slab preset's
% own showcase size 4x4 sits above the drift formula (3 x 0.60 = 1.8 >
% 1.5), the width cost the sweep priced into the preset; the guard
% honestly reports it and the picture is drawn regardless.
\documentclass[varwidth=20cm, border=6pt]{standalone}
\usepackage{amsmath, amssymb}
\usepackage{tenkz}
\begin{document}

% ---- item 1: derived sheet sep + planes-local ratios -------------------
% (1) 3x3 double layer, NO sheet sep: the derived default
% sep = (rows-1) x planerise + planegap = 2 x 0.45 + 0.70 = 1.60 must
% reproduce the z_R1-class geometry (no mid-band dot near-merge; the old
% static 1.4 left less than daylight between the default beads).
% The planes-local ratios are slant 0.40 / rise 0.45 (leg-to-dot
% clearance 0.2 lp for every pairing leg climbing up to 4 row-levels).
\textbf{(1) planes, derived sep, 3x3}\par\smallskip
\begin{tenkzplanes}[rows=3, cols=3, sheets={ket, bra}]
\end{tenkzplanes}

\medskip
% (2) 4x4 with an open column at the derived sep: the sheets must NOT
% interleave (the z_C1 regression: static 1.4 < 3 x rise).  Derived
% sep = 3 x 0.45 + 0.70 = 2.05.
\textbf{(2) planes, derived sep, 4x4, open column 2}\par\smallskip
\begin{tenkzplanes}[rows=4, cols=4, sheets={ket, bra}, open={(1-4,2)}]
\end{tenkzplanes}

\medskip
% (3) explicit sheet sep=1.65 at 5x5: below (rows-1) x rise = 4 x 0.45
% = 1.80, so the interleave guard must WARN (the z_X2 catastrophe: sep at
% an integer multiple of the rise leaves less than daylight between the
% audited outer bead inks) yet still draw what it was told.
\textbf{(3) planes, sheet sep=1.65 at 5x5 (warns: interleave + sheet-coincide)}%
\par\smallskip
\begin{tenkzplanes}[rows=5, cols=5, sheets={ket, bra}, sheet sep=1.65]
\end{tenkzplanes}

\medskip
% ---- item 2: presets, overrides, lean ---------------------------------
% (4) default oblique sheet at the retuned rise 0.60: physical-leg tips
% must end >= 0.09 lp short of the behind-row bond (rise must exceed
% physleg/latticepitch = 0.38/0.75 ~ 0.507; the old 0.55 tips ended
% 0.043 lp short of the bond -- a near-touch, plane_sweep1 zcpt_55).
\textbf{(4) oblique default (slant 0.45 / rise 0.60)}\par\smallskip
\begin{tenkzlattice}[rows=4, cols=4, frame=oblique, physical=up,
                     boundary legs]
\end{tenkzlattice}

\medskip
% (5) frame=slab (slant 0.60 / rise 0.55): the full-pitch RMP showcase
% geometry (plane_sweep1 zsw_6055 -- 47-degree leg/bond separation), with
% a filled region to prove the tracer shears with the steeper frame.
\textbf{(5) frame=slab with region $R$}\par\smallskip
\begin{tenkzlattice}[rows=4, cols=4, frame=slab, physical=up,
                     boundary legs]
  \tnregion[slot=selected, label=$R$]{(2-3, 2-3)}
\end{tenkzlattice}

\medskip
% (6) explicit overrides OUTSIDE the sane ranges: plane rise=0.40 is at
% or below the 0.507 leg threshold (worst sweep cell 0.30/0.40: false
% 5-valent junctions) and plane slant=0.25 is below the readable band
% (depth bonds steepen to 61-67 degrees).  Both keys must WARN yet draw.
\textbf{(6) plane slant=0.25 / plane rise=0.40 (warns: both guards)}%
\par\smallskip
\begin{tenkzlattice}[rows=3, cols=4, frame=oblique, physical=up,
                     plane slant=0.25, plane rise=0.40]
\end{tenkzlattice}

\medskip
% (7) plane lean=west: the mirrored "\" sheet.  The frame map stays
% linear under the sign flip, so the region tracer's fills, rounded
% corners, margins and outside-corner label all survive (plane_sweep1
% page 2a, zmir_bare/zmir_reg).
\textbf{(7) plane lean=west with region $R$}\par\smallskip
\begin{tenkzlattice}[rows=4, cols=4, frame=oblique, plane lean=west,
                     physical=up, boundary legs]
  \tnregion[slot=selected, label=$R$]{(2-3, 2-3)}
\end{tenkzlattice}

\medskip
% ---- item 3: tall-window guard ----------------------------------------
% (8) 5x3 at stock slant: (rows-1) x slant = 4 x 0.45 = 1.80 exceeds
% 0.5 x (cols-1) = 1.0, so the window collapses toward a rhombus
% staircase (plane_sweep1 page 2b, drift/width 0.9) and must WARN.
\textbf{(8) 5x3 tall window (warns: rhombus collapse)}\par\smallskip
\begin{tenkzlattice}[rows=5, cols=3, frame=oblique]
\end{tenkzlattice}

\medskip
% ---- items 4 + 5: frame=vertical, trace style ------------------
% (9) the RMP O_L word as a ONE-LINER (arXiv:2011.12127 Fig:MPDO-O_L):
% vertical frame (columns descend, \tndots renders \vdots), op-row boxes
% whose virtual legs point west/east, and the ring closure as truncated
% half-hooks -- hooks is the vertical-frame DEFAULT for periodic (a
% vertical racetrack would spend the width the rotation saves), so
% `periodic` alone suffices; `periodic, trace style=hooks` is spelled out here
% to exercise the key itself.
\textbf{(9) vertical $O_L$ word, hooks closure}\par\smallskip
\[
  O_L(M) \;=\;
  \begin{tenkz}[frame=vertical, physical=updown,
                periodic, trace style=hooks, tensor style=box]
    \tn{M} & \tndots & \tn{M}
  \end{tenkz}
\]

\medskip
% (10) the rotated transfer-matrix column E (RMP sec2fig8, plane_sweep3
% panel 2b): a ket-over-bra sandwich in the vertical frame -- the two
% dots sit side by side, their physical pairing runs horizontally --
% via the `chain axis=south` alias spelling of the frame key.
\textbf{(10) rotated transfer matrix, chain axis=south}\par\smallskip
\[
  E \;=\;
  \begin{tenkz}[chain axis=south, sandwich]
    \tn{A}\\
    \tn{A}
  \end{tenkz}
\]

\medskip
% (11) periodic, trace style=racetrack spelled explicitly on a HORIZONTAL word:
% the established return-wire closure must be reachable by name (it is
% the horizontal-frame default; the key also selects periodic closure).
\textbf{(11) horizontal word, periodic, trace style=racetrack}\par\smallskip
\[
  \operatorname{Tr}\,ABC \;=\;
  \begin{tenkz}[periodic, trace style=racetrack]
    \tn{A} & \tn{B} & \tn{C}
  \end{tenkz}
\]

\medskip
% ---- item 6: per-side boundary policy ---------------------------------
% (12) the hard11 brick-circuit fragment with west=none: prepared /
% discarded indices are SIDES of the word, not rows -- the fragment
% seals its west (input) indices and keeps its east (output) indices
% open, so NO west stubs render (hard11_v3's four spurious west stubs
% gone) while all four east stubs stay.
\textbf{(12) brick circuit, west=none / east=open}\par\smallskip
\[
  |\psi\rangle \;=\;
  \begin{tenkz}[rows={wire,wire,wire,wire}, west=none, east=open]
     & & \tn[box, wires=2]{}\\
    \tn[dot]{|0\rangle} & \tn[box, wires=2]{} & \\
     & & \tn[box, wires=2]{}\\
     & &
  \end{tenkz}
\]

\medskip
% (13) the converse split, west=open / east=none, on a plain word: the
% side policies must resolve independently per side (a row whose east is
% `none` keeps its west stub).
\textbf{(13) plain word, west=open / east=none}\par\smallskip
\[
  \begin{tenkz}[west=open, east=none]
    \tn{A} & \tn{B}
  \end{tenkz}
\]

\medskip
% ---- item 7: free-tier label pos --------------------------------------
% (14) label pos on \tnput: the rotated-E mock of plane_sweep3 panel 2a
% -- a vertical wire strikes the historical south band at every bead, so
% both labels yield east (no strike-through); plus three free-standing
% beads, one per remaining compass value (default south unchanged,
% north, west).
\textbf{(14) free tier, label pos}\par\smallskip
\begin{tenkzfree}
  \tnput[label pos=east]{K}{(0,0)}{A}
  \tnput[label pos=east]{B}{(0,-1.1)}{\bar A}
  \tnjoin{K.south}{B.north}
  \tnput{S}{(1.4,-0.55)}{S}
  \tnput[label pos=north]{N}{(2.4,-0.55)}{N}
  \tnput[label pos=west]{W}{(3.4,-0.55)}{W}
\end{tenkzfree}

\end{document}
